
\section{\sys{} Implementation}

In this section, we describe the key implementation details of \sys, including kernel modifications, GPU runtime implementation, and interaction with existing GPU frameworks.
\label{sec:implementation}

\subsection{Kernel Modifications}

\sys requires minimal but carefully placed modifications to \sysdriver to expose programmable hook points without compromising system stability or performance.

\sys leverages the NVIDIA Open GPU Kernel Modules, which have become the recommended default for modern datacenter GPUs (Turing, Ampere, Hopper, and newer). By building upon this open-source foundation, \sys avoids brittle binary patching of proprietary drivers. Instead, it introduces standard, maintainable extension points within the open source driver code, aligning with the vendor's shift towards open kernel interfaces while maintaining compatibility with the closed-source user-space CUDA runtime.

\paragraph{Specific hooks and patches:}
We add specific hooks within the NVIDIA UVM module at critical memory management decision points, including page fault handlers, prefetch logic, and eviction paths. Similarly, we extend kernel-scheduler entry points to enable custom admission control and resource allocation policies. These modifications are designed to be minimal, non-intrusive, and compatible with existing driver functionality.

All hook points follow the struct\_ops-style callback pattern, where verified eBPF programs can be dynamically attached and detached without requiring driver recompilation or system restarts. The hooks provide access to relevant context (e.g., fault address, process ID, memory pressure metrics) while maintaining safety through the eBPF verifier.

\subsection{GPU Runtime Implementation Details}

\paragraph{PTX JIT compiler integration:}
\sys integrates a PTX Just-In-Time (JIT) compiler that translates verified eBPF bytecode into optimized PTX instructions. PTX (Parallel Thread Execution) is an intermediate representation for NVIDIA GPUs that provides better performance than binary-level instrumentation frameworks. The JIT compiler produces architecture-specific code optimized for the target GPU, enabling dynamic instrumentation with minimal overhead.

\paragraph{Runtime resource bounding:}
The GPU runtime enforces strict resource bounds to prevent injected code from interfering with application kernels. This includes limiting instruction counts per invocation, bounding register usage to avoid spilling, and restricting memory access to designated shared regions. These bounds are verified at load time and enforced at runtime through hardware mechanisms and software checks.

\paragraph{Interaction with existing GPU frameworks:}
\sys is designed to be compatible with existing GPU programming frameworks and tools. It works seamlessly with CUDA applications and can coexist with CUPTI and NVBit instrumentation when needed. The trampoline mechanism and shared memory infrastructure are transparent to application code, requiring no modifications to existing GPU kernels.

\subsection{Shared Memory}

We place maps and coordination buffers in \texttt{managed\_shared\_memory} for zero-copy CPU/GPU access, reducing marshaling and latency.


\subsection{Cross-Domain Memory Synchronization}

Cross-domain visibility requires explicit barriers (device \_\_threadfence\_system and host stream sync). To transfer writes from the GPU to the CPU, the device must execute \_\_threadfence\_system() to flush any outstanding writes into system-visible memory, followed by a host-side cudaStreamSynchronize() to wait for the GPU to complete. We encapsulate these as a simple device/host pair to centralize tuning and reduce boilerplate:
\begin{minted}[highlightlines=false,fontsize=\small]{c}
// Ensure GPU writes become visible to the CPU
__device__ inline void memsys_bar_gpu() {
    __threadfence_system();
}
// Wait for the GPU on the given stream
static inline void memsys_bar_cpu(cudaStream_t stream =
        0) {
    cudaStreamSynchronize(stream);
}
\end{minted}
In a persistent kernel, call the GPU fence before reading shared state; after host updates, synchronize the stream to ensure visibility. This two-call barrier hides cross-domain complexity. On the host side, memsys\_bar\_cpu() calls cudaStreamSynchronize() on a specified stream to wait for the GPU to finish its pending commands, including memory writes. This wrapper simplifies code, clarifies intent, and centralizes performance tuning.
\begin{minted}[highlightlines=false,fontsize=\small]{c}
using namespace boost::interprocess;
__device__ int *shared_flag;
__global__ void persistent_kernel() {
    memsys_bar_gpu();// Ensure CPU writes are visible
    int cmd = *shared_flag;
    // Jump into PTX injector...
}
int main() {
    managed_shared_memory shm(open_or_create, "MyShm"
        , 1 << 20);
    int *flag = shm.construct<int>("Flag")(0)[1];
    *flag = 42;              // Write command
    memsys_bar_cpu(0);       // Ensure GPU sees update
    persistent_kernel<<<grid, block>>>();
    cudaDeviceSynchronize();
}
\end{minted}

\subsection{Synchronization}
We use host atomics/interprocess primitives and device CUDA atomics/fences over well-defined layouts in shared memory. To reduce contention, only a subset of threads performs atomics while others proceed in parallel.

\subsection{PTX Generation \& Injection}\label{sec:ptx_injection}

User/eBPF logic triggers JIT generation of PTX snippets that are installed via launch-time patching ("PTX injection"). Snippets interoperate with eBPF maps and shared buffers to support flexible, low-latency instrumentation without restarting the application.

\paragraph{Trampoline Mechanism}
At runtime we patch candidate sites (entry, memory access, barriers) with a tiny PTX stub that:
\begin{itemize}[leftmargin=*,itemsep=0pt]
  \item Saves the minimal GPU register and predicate state.
  \item Jumps into the injected instrumentation snippet.
  \item Restores the original state and resumes normal execution.
\end{itemize}
These trampolines are linked into the device's code image via launch-time patching of branch targets, enabling fine-grained hooking without requiring application restart or full recompilation. Our trampoline generator handles label resolution, register allocation, and insertion of necessary synchronization primitives to avoid hazards.



\subsection{Isolation Requirement}\label{sec:isolation}

Launch-time PTX injection and trampoline-based instrumentation introduce new code into GPU kernels, which might bring the risk of corrupting application state or violating memory safety. \sys therefore runs injected snippets in a protected domain.

\paragraph{Runtime Sandbox Integration}
During JIT compilation, \sys{} wraps each PTX snippet in a minimal prologue/epilogue sequence that:
%\begin{itemize}[leftmargin=*,itemsep=0pt]
  %\item
  a) Saves a live register state to a safe context.
  %\item
  b) Applies bitwise partition masking and invokes the instrumentation logic.
  %\item
  c) Restores registers and resumes normal kernel execution.
%\end{itemize}

Combining trampolines and shared-memory coordination ensures injected code remains isolated from application kernels.
